\section{A repaired base-change step in the symplectic construction}
\label{sec:s5-symplectic-source-repair}

\paragraph{Literal source finding.}
Calegari's Lemma~4.2, author-hosted publisher PDF pp.~9--10,
contains on p.~9 the unnumbered equality
$\wedge^2\Pi_E=\mu\boxplus\mu$ in its hypothetical
cuspidal-induction branch, followed by a multiplicity-two
argument.  The same passage is in the supplied
\texttt{S5ArtinWeb.tex}, lines~563--572, under
\texttt{lemma:satake}.  The published page was checked
both as text and as a rendered image in the present audit.
The correct quadratic base change is
$\mu\boxplus\mu^\sigma$, where $\sigma$ is the nontrivial
element of $\operatorname{Gal}(E/\mathbb Q)$.
The identical-summand equality and its stated local-multiplicity
justification are not valid.  The intended conclusion has the
following complete repair using the same classical inputs.

\paragraph{MA-S5-008 (corrected summands and a constituent-degree repair).}
In Calegari's specified family, let $\pi_E$ be the established
cuspidal representation of $\operatorname{GL}_2(\mathbb A_E)$,
put $\omega=\omega_{\pi_E}$, and retain the established cuspidal
descent
\[
 \operatorname{BC}_{E/\mathbb Q}\Pi
        =\operatorname{Sym}^3\pi_E.
 \tag{MA-S5.40}
\]
Its exterior-square transfer $\Sigma=\wedge^2\Pi$ is
noncuspidal.  In fact it has the exact isobaric type
\[
 \Sigma=\chi\boxplus\Sigma_5,
 \qquad \Sigma_5\text{ cuspidal of degree five},
 \tag{MA-S5.41}
\]
with
\[
 \operatorname{BC}\chi=\omega^3,
 \qquad
 \operatorname{BC}\Sigma_5=\operatorname{Sym}^4\pi_E\otimes\omega.
 \tag{MA-S5.42}
\]
The representation $\Pi$ is of symplectic type, with similitude
character equal to the character $\nu$ fixed immediately before
the source's Lemma~4.2 by $\Pi^\vee\simeq\Pi\otimes\nu^{-1}$.
Consequently its unramified Satake parameters have the asserted
form
\[
                         \{a,b,\nu(p)/a,\nu(p)/b\}.
 \tag{MA-S5.43}
\]
This repair uses no comparison with $W$ at the exceptional
primes and no zero-location or real nonvanishing hypothesis.
In particular it does not use the later conclusion
$\nu=\psi^{-1}$ of Calegari's Lemma~4.6.

\begin{proof}
\emph{The actual base change has constituent degrees one and five.}
For a two-dimensional vector space representation $\tau$,
the polynomial identity is
\[
 \wedge^2\operatorname{Sym}^3\tau
   =(\operatorname{Sym}^4\tau\otimes\det\tau)
                                      \oplus(\det\tau)^3.
 \tag{MA-S5.44}
\]
For example, if $a,b$ are the two eigenvalues, the six
exterior products of $a^3,a^2b,ab^2,b^3$ are
\[
                 a^5b,\ a^4b^2,\ a^3b^3,\ a^3b^3,
                                  \ a^2b^4,\ ab^5.
 \tag{MA-S5.45}
\]
Five are the weights of $(ab)\operatorname{Sym}^4\tau$;
the remaining one is $(ab)^3$.  Equality on the diagonal
torus proves (MA-S5.44) as a semisimple polynomial
representation identity.

Kim's exterior-square transfer and the already used
symmetric-fourth transfer exist as isobaric automorphic
representations.  In this icosahedral case
$\operatorname{Sym}^4\pi_E$ is cuspidal, as in the opening
of Calegari's \S4.  Cyclic base change and exterior powers
commute at every common unramified place: both apply the
same exterior-power map to the local parameter restricted
to $E$.  Thus isobaric strong multiplicity one and
(MA-S5.40)--(MA-S5.45) give
\[
 \operatorname{BC}\Sigma
    =\omega^3\boxplus
                       (\operatorname{Sym}^4\pi_E\otimes\omega).
 \tag{MA-S5.46}
\]
This is an isobaric decomposition into a character and a
cuspidal degree-five representation.  It is not a claim
about double eigenvalue multiplicities at each prime.

Suppose first that $\Sigma$ were cuspidal of degree six.
Its noncuspidal quadratic base change forces
\[
 \Sigma=\operatorname{AI}_{E/\mathbb Q}\mu,
 \qquad
 \operatorname{BC}\Sigma=\mu\boxplus\mu^\sigma,
 \tag{MA-S5.47}
\]
where $\mu$ is cuspidal of degree three.  The second
identity is the restriction-of-induction formula at
unramified parameters, promoted by isobaric multiplicity
one; the two cosets give $\mu$ and its Galois conjugate.
Both constituents have degree three.  They cannot equal
the degree-one and degree-five cuspidal constituents in
(MA-S5.46), by uniqueness of isobaric decomposition.
Hence $\Sigma$ is noncuspidal.

Moreover $\mu\not\simeq\mu^\sigma$ in the hypothetical
cuspidal induction in (MA-S5.47): a Galois-invariant
cuspidal representation descends, and its induction is
the sum of the two quadratic twists of a descent.
Thus in precisely that branch the source cannot replace
the conjugate summand by an identical one.  This explains
the literal error independently of the repaired
constituent-degree contradiction.

To prove the stronger assertion (MA-S5.41), decompose
$\Sigma$ into its cuspidal isobaric constituents over
$\mathbb Q$.  Under quadratic base change a cuspidal
constituent of degree $r$ either remains cuspidal of
degree $r$, or splits into two cuspidal constituents of
degree $r/2$.  In order to contribute the degree-five
constituent of (MA-S5.46), a constituent over $\mathbb Q$
must have degree five or ten.  The total degree is six,
so it has degree five.  The remaining degree is one.
This proves (MA-S5.41), and uniqueness of the two
constituents after base change gives (MA-S5.42).

\emph{The classification now selects symplectic type explicitly.}
Apply the same published cuspidality classification that
the source invokes: Asgari--Raghuram, Theorem~1.1(ii).
We check its alternatives using its proved exterior-square
decompositions, rather than assuming that every noncuspidal
exterior square is symplectic.

In its tensor-product case (a), Proposition~3.1(a) gives
two isobaric representations of degree three.  None of
their cuspidal constituents has degree five.
In its dihedral Asai case (b), Proposition~3.2 gives
the automorphic induction of a twisted symmetric square
of a dihedral degree-two representation.  That symmetric
square is an isobaric sum of degrees one and two, possibly
further split.  Quadratic automorphic induction therefore
has constituents of degree at most four.
Finally, in its quadratic-induction case (d), Proposition~3.5
gives an Asai representation of degree four together with
an induced character of degree two.  Again every cuspidal
constituent has degree at most four.
All three alternatives contradict the cuspidal degree-five
constituent in (MA-S5.41).

The remaining alternative (c) says that $\Pi$ is a
transfer from a globally generic cuspidal representation
of $\operatorname{GSp}_4$.  If its similitude character
is $\widetilde\nu$, Proposition~3.3 gives
$\wedge^2\Pi=\widetilde\nu\boxplus\widetilde r_5$.
The unique degree-one constituent of (MA-S5.41) is
therefore $\widetilde\nu=\chi$; the remaining constituent
is $\Sigma_5$.  In particular
\[
                 \Pi^\vee\simeq\Pi\otimes\chi^{-1},
                 \qquad \omega_\Pi=\chi^2.
 \tag{MA-S5.48}
\]

\emph{Identification of the previously fixed similitude.}
The same constituent argument excludes quadratic
automorphic induction of $\Pi$ from \emph{any} quadratic
field, since Proposition~3.5 would always give only
constituents of degree at most four in $\wedge^2\Pi$.
Suppose $\Pi\simeq\Pi\otimes\xi$ for a nontrivial
Hecke character $\xi$.  Comparing central characters
gives $\xi^4=1$.  Either $\xi$ is quadratic, or
$\xi^2$ is a nontrivial quadratic self-twist.
The cyclic self-twist criterion, also proved as
Asgari--Raghuram Lemma~5.2, would then make $\Pi$ a
quadratic automorphic induction, a contradiction.
Thus $\Pi$ has no nontrivial character self-twist.

Compare (MA-S5.48) with the already established
$\Pi^\vee\simeq\Pi\otimes\nu^{-1}$.
It gives $\Pi\simeq\Pi\otimes\chi\nu^{-1}$, so
the absence of self-twists proves $\chi=\nu$.
This supplies the last step in the original lemma
without its unexplained quadratic-induction shortcut.

For clarity, the asserted parameter form holds at every
unramified place of $\Pi$, not only at places used in a
weak transfer.  Local self-duality with multiplier $\nu$
first makes $\nu$ unramified there and makes the Satake
multiset invariant under $z\mapsto\nu(p)/z$.
Its product is $\nu(p)^2$ by (MA-S5.48).
Nonfixed orbits already form pairs with product $\nu(p)$.
The remaining values are $\sqrt{\nu(p)}$ and
$-\sqrt{\nu(p)}$.  The product condition makes the
multiplicity of the negative value even; since the
remaining cardinality is even, the positive multiplicity
is even too.  Pair equal fixed values.  This proves
(MA-S5.43) in all cases, including repeated eigenvalues.
\end{proof}

\paragraph{Role, status, and exact source scope.}
The corrected conclusion is the symplectic parameter input
used in MA-S5-003/006 and MA-RB-005.  It restores that
input while leaving the exceptional-prime comparison and
the shared-real-zero gate unchanged.  No claim is made
that the earlier source readings detected this error:
the literal summand mistake was identified in the new
read for MA-RB-005 and audited separately here.
The outcome is a refuted printed intermediate equality
with a proved refinement, not a refuted construction theorem.

The exact primary Calegari locator is Lemma~4.2, publisher
PDF pp.~9--10, especially the unnumbered base-change
equality at the foot of p.~9.  The required symmetric-cube
descent is on p.~9 before Lemma~4.1, and symmetric-fourth
cuspidality is stated at the opening of \S4, p.~8.
The supplied source's plethysm also occurs in the proof
of Lemma~2.5; (MA-S5.44)--(MA-S5.45) reprove it here.
The author-hosted publisher version is
\url{https://math.uchicago.edu/~fcale/papers/S5.pdf}.
The source's cyclic base-change input remains
Arthur--Clozel, Theorem~4.2, printed p.~202;
its full trace-formula proof is not newly audited.

For the completion of the repair, the primary
Asgari--Raghuram author version
\url{https://math.okstate.edu/people/asgari/Files/cuspidality.pdf}
was freshly read at Theorem~1.1, p.~2; isobaric
multiplicity one, Theorem~2.1, p.~4; the complete proofs
of Propositions~3.1--3.3 and~3.5, pp.~8--10; and the
complete self-twist proof, Lemma~5.2, p.~11.
The main cuspidality classification itself, Kim's
transfer, symmetric-power cuspidality, and cyclic base
change are the same carried published inputs as in the
source construction.  They are not presented as new
unconditional functoriality proved by this audit.

Calegari, F. (2013). The Artin conjecture for some
$S_5$-extensions. \emph{Mathematische Annalen, 356},
191--207. \url{https://doi.org/10.1007/s00208-012-0839-4}.

Asgari, M., \& Raghuram, A. (2011).
A cuspidality criterion for the exterior square transfer
of cusp forms on $\mathrm{GL}(4)$.
In \emph{On certain $L$-functions} (Clay Mathematics
Proceedings, Vol.~13, pp.~33--53).
American Mathematical Society.

Arthur, J., \& Clozel, L. (1989).
\emph{Simple algebras, base change, and the advanced
theory of the trace formula}.
Annals of Mathematics Studies, Vol.~120.
Princeton University Press.
